%%%
%
% $Autor: Wings $
% $Datum: 2021-05-14 $
% $Pfad: ArduinoIDE2 $
% $Dateiname: 
% $Version: 4620 $
%
% !TeX spellcheck = de_DE/GB
% !TeX program = pdflatex
% !BIB program = biber/bibtex
% !TeX encoding = utf8
%
%%%


\chapter{Arduino IDE 2.3.6}\label{ArduinoIDE2}
Dieses Kapitel bietet eine umfassende Untersuchung der Arduino IDE, einschließlich ihrer Funktionen, Schnittstellenoptionen und Funktionalität. Es enthält eine detaillierte Erklärung des Zwecks und der Nutzbarkeit jeder Komponente, eine schrittweise Installationsanleitung sowie Anweisungen zur Initialisierung und Konfiguration der Umgebung. Dieser grundlegende Überblick befähigt die Leser, die IDE effektiv für die Programmierung und Fehlerbehebung von Arduino-kompatibler Hardware zu nutzen.

\section{Beschreibung der Arduino IDE}\label{ArduinoIDE}

Die Arduino IDE ist eine Open-Source-Software, die speziell für das Bearbeiten, Kompilieren und Hochladen von Code auf Arduino-Module entwickelt wurde. Sie ist plattformübergreifend und mit Betriebssystemen wie Windows, Linux und macOS kompatibel. Auf der Java-Plattform aufbauend unterstützt die IDE verschiedene Arduino-Module und die Programmierung in C und C++.

Die Mikrocontroller auf den Arduino-Boards werden so programmiert, dass sie die in Form von Code bereitgestellten Informationen verarbeiten. Programme, die in der IDE geschrieben werden, werden als Skizzen (Sketches) bezeichnet. Diese Skizzen generieren eine Hex-Datei, die dann auf den Mikrocontroller übertragen und hochgeladen wird.

Die IDE-Umgebung besteht aus zwei Hauptkomponenten: dem Editor und dem Compiler. Der Editor wird zum Schreiben von Code verwendet, während der Compiler den Code kompiliert und auf das Arduino-Modul hochlädt.\cite{Fezari:2018} In Version 2.3.3 der Arduino IDE spielt die Menüleiste, die sich am oberen Rand der Schnittstelle befindet, eine entscheidende Rolle, da sie Zugriff auf wichtige Befehle wie Hochladen, Verifizieren/Kompilieren und Speichern bietet. Zusätzlich enthält sie das Dateimenü zum Öffnen neuer oder bestehender Dateien sowie den Beispiele-Bereich, der vorgefertigte Skizzen für verschiedene Anwendungen wie Blink und Fade bereitstellt.

\bigskip
Die 6 Buttons am oberen Rand des Fensters sind wie folgt in \ref{fig:Menü-Ansicht}:

\begin{figure}[H]
	\centering
	\caption{Menü-Leiste}
	\label{fig:Menü-Ansicht}
	\includegraphics[width=\textwidth]{Arduino/ArduinoIDE/menuAnzeige}
\end{figure}
		


\bigskip

\begin{itemize}
	\item[\hbox{
\includegraphics[viewport=6 0 34 30, clip=true,height=12pt]{Arduino/ArdiunoIDE2/MenuButton}}] Das Häkchen-Symbol dient zum \menu{Überprüfen} des geschriebenen Codes. Nach der Code-Eingabe prüft diese Funktion auf Fehler und bestätigt die korrekte Struktur. Dieser Schritt stellt sicher, dass der Code für die Hardware bereit ist.
	
	\item[\hbox{
\includegraphics[viewport=34 0 63 30, clip=true,height=12pt]{Arduino/ArdiunoIDE2/MenuButton}}] Das \menu{Hochladen}-Symbol kompiliert den Code und überträgt ihn auf das angeschlossene Board, sodass dieser ausgeführt werden kann.
	
	\item[\hbox{
\includegraphics[viewport=60 0 88 30, clip=true,height=12pt]{Arduino/ArdiunoIDE2/MenuButton}}] Das \menu{Debug}-Symbol startet eine Debugging-Session, mit der Sie den Code analysieren können – inklusive Breakpoints, Schritt-für-Schritt-Ausführung und Variableninspektion (auf kompatiblen Boards).
	
	\item[\hbox{
\includegraphics[viewport=702 0 730 30, clip=true,height=12pt]{Arduino/ArdiunoIDE2/MenuButton}}] Das \menu{Serieller Plotter}-Symbol öffnet ein Tool zur Visualisierung von Echtzeitdaten, die über die serielle Verbindung vom Board gesendet werden, und stellt sie als Graphen dar.
	
	\item[\hbox{
\includegraphics[viewport=720 0 800 30, clip=true,height=12pt]{Arduino/ArdiunoIDE2/MenuButton}}] Das \menu{Serieller Monitor}-Symbol öffnet ein Terminal zur Anzeige und zum Versand von Textdaten über die serielle Verbindung, ermöglicht also die Echtzeit-Kommunikation mit dem Board.
\end{itemize}


\section{Installation der Arduino IDE}
\subsection{Download der Arduino IDE}
Der Arduino Nano 33 BLE Sense nutzt die Arduino Software Integrated Development Environment (IDE) für die Programmierung, die am weitesten verbreitete IDE für alle Arduino-Boards und sowohl online als auch offline verwendbar ist.

Diese Open-Source-Arduino-IDE vereinfacht das Schreiben von Code und das Hochladen auf das Board. Für jedes Betriebssystem (OS), einschließlich macOS, Linux und Windows, sind verschiedene Versionen der Software verfügbar. Die Arduino-Community bietet zudem eine Online-Plattform zum Programmieren und Speichern von Sketches in der Cloud. Dieser Online-Arduino-Editor ist die neueste Version der IDE und beinhaltet alle Bibliotheken sowie Unterstützung für neue Arduino-Boards.

Um auf diese Softwarepakete zuzugreifen, besuchen Sie die folgende Website:
\HREF{https://www.arduino.cc/en/software}{Arduino-Software},
um auf dem neuesten Stand zu bleiben, da tägliche Updates über den genannten Link verfügbar sind.

Der Editor kann direkt von der Arduino-Software-Seite heruntergeladen werden. Die Download-Optionen sind in Abbildung ~\ref{fig:downloadArduinoIDE} dargestellt.

\begin{figure}[H]
	\centering
	\caption{Arduino IDE 2.3.6}
	\label{fig:downloadArduinoIDE}
	\includegraphics[width=\textwidth]{Arduino/ArduinoIDE/downloadArduinoIDE.png}
\end{figure}


%\subsection{Download-Optionen}
%Below is a summary ~\ref{fig:Summary}  of the available download options for the Arduino IDE, tailored to different operating systems and individual needs. Choose the appropriate version for the device to ensure compatibility and access to the latest features.

%\begin{center}\centering
%	\includegraphics[width=12cm]{Arduino/ArdiunoIDE2/SummaryOfOptions.png}
%	\captionof{figure}{Summary of DownloadOptions}
%	\label{fig:Summary}		
%\end{center}

\subsection{Installation auf Windows}
Der Installationsprozess der Arduino IDE 2 auf einem Windows-Computer wird im folgenden Abschnitt Schritt für Schritt beschrieben. 
So installieren Sie die Arduino IDE 2 auf einem Windows-Computer:

Führen Sie die von der Software-Seite \HREF{https://www.arduino.cc/en/software}{Arduino-Software} heruntergeladene Datei aus.

Nutzen Sie die Installationsmöglichkeiten unter Windows, wie in Abbildung ~\ref{fig:downloadArduinoIDErot} dargestellt, um eine korrekte Installation sicherzustellen.

\begin{figure}[H]
	\centering
	\caption{Option für Windows}
	\label{fig:downloadArduinoIDErot}
	\includegraphics[width=\textwidth]{Arduino/ArduinoIDE/downloadArduinoIDErot.png}
\end{figure}

Nachdem der Download abgeschlossen ist, führen Sie die Setup-Datei aus. Akzeptieren Sie das Lizenzabkommen \ref{fig:lizenzabkommen} und klicken sie auf \menu{Annehmen}:

\begin{figure}[H]
	\centering
	\caption{Lizenzabkommen}
	\label{fig:lizenzabkommen}
	\includegraphics[width=\textwidth]{Arduino/ArduinoIDE/lizenzabkommen.png}
\end{figure}

Danach wählen Sie aus, für wen die Installation auf Ihrem PC angelegt werden soll, Sie können auswählen, ob für Sie selbst oder alle Benutzer des Computers: \ref{fig:downloadBenutzer}

\begin{figure}[H]
	\centering
	\caption{Auswahl der Benutzer}
	\label{fig:downloadBenutzer}
	\includegraphics[width=\textwidth]{Arduino/ArduinoIDE/downloadBenutzer.png}
\end{figure}

Nachdem Sie Ihre Wahl getroffen haben, klicken Sie auf \menu{Weiter}. Nun können Sie das Verzeichnis auswählen, in das die Arduino IDE installiert werden soll. \ref{fig:downloadVerzeichnis}

\begin{figure}[H]
	\centering
	\caption{Auswahl des Installationsverzeichnis}
	\label{fig:downloadVerzeichnis}
	\includegraphics[width=\textwidth]{Arduino/ArduinoIDE/downloadVerzeichnis.png}
\end{figure}

Klicken Sie dann auf \menu{Installieren} und warten Sie ab, bis der Download abgeschlossen ist. Nach Abschluss des Downloads öffnet sich die Arduino IDE automatisch. \ref{fig:downloadIDE}

\begin{figure}[H]
	\centering
	\caption{Die Arduino IDE}
	\label{fig:downloadIDE}
	\includegraphics[width=\textwidth]{Arduino/ArduinoIDE/downloadIDE.png}
\end{figure}


\subsection{Installation (MacOS)}
Aufgrund fehlender Hardware können wir leider keine Installation auf MacOS bereitstellen.

%To install the Arduino IDE, begin by downloading the latest version from the Arduino website Arduino Software \HREF{https://www.arduino.cc/en/software}{Arduino-Software}. Select the version compatible with the specific operating system in use. In this case, Arduino 2.3.2 is being installed for macOS (Sonoma 14.4.1). The setup file is named \HREF{https://www.arduino.cc/en/software}{Arduino-Software} and has a size of 193,600 KB. This file is in Zip format. For those downloading Arduino IDE 2.3.X with Safari, the file will automatically extract upon download completion, while other browsers may require manual extraction. The figure below displays the latest offline version of Arduino IDE 2.3.3 ~\ref{fig:ArduinoIDECreateAgentInstallation}, which is also compatible with all operating systems.




%\begin{center}
	
	%\includegraphics[width=\linewidth]{Arduino/ArdiunoIDE2/macDownload.png}
	%\label{fig:ArduinoIDE Create Agent Installation}
	%\captionof{figure}{ArduinoIDE Create Agent Installation}
	%\label{fig:ArduinoIDECreateAgentInstallation}
	
	
	
	
	
	
	%\Mynote{Die Installation muss aktualisiert werden, die folgenden Bilder müssen mit einem Mac ebenso aktualisiert werden}
	
%\end{center}

%\begin{center}
%	\includegraphics[width=0.7\linewidth]{Arduino/ArdiunoIDE2/OpenTheDownloadFolder.png}
%	\captionof{figure}{Open the Download folder}
%\end{center}

%Copy the Arduino application bundle into the application's folder (or elsewhere on the
%computer) then it looks like Figure   ~\ref{fig:CopytotheAppliationsFolder}.

%\begin{center}
	
%	\includegraphics[width=0.7\linewidth]{Arduino/ArdiunoIDE2/CopyToTheApplicationsFolder.png}
%	\captionof{figure}{Copy to the Applications Folder}
%	\label{fig:CopytotheAppliationsFolder}
%\end{center}

%It can be seen from the Figure ~\ref{fig:ArduinoIDESketch2} that the basic Arduino sketch has two parts. 

%\begin{itemize}
	
%	\item \PYTHON{void setup()}: This function returns void
%	and performs initializations such as setting the output LED color and specifying the core.
%	\item \PYTHON{void loop()}: In this function, specific operations to be executed within the loop are defined. The code for each operation is enclosed within curly braces \PYTHON{\{\}}, and each function has a return type. In this case, the return type is \menu{void}.
	
%\end{itemize}

%\begin{center}
	
	
	
%	\includegraphics[width=0.7\linewidth]{Arduino/ArdiunoIDE2/ArduinoIDESketch2.png}
%	\captionof{figure}{Arduino IDE Sketch}
%	\label{fig:ArduinoIDESketch2}
	
	
	
%	\Mynote{Figure 1.22 muss größer und alle Bilder in diesem Kapitel ohne Schwarzmodus}	
	
%\end{center}





\section{Übersicht}
Die Arduino IDE 2 beeinhaltet eine neue Sidebar \ref{fig:ArduinoSidebar}, auf der die am häufigsten genutzten Werkzeuge leichter zu erreichen sind. Folgende Werkzeuge sind besonders hilfreich: \cite{ArduinoWeb:2025}

\begin{itemize}
	\item \textbf{Überprüfen / Hochladen} Diese Funktionen dienen zum Kompilieren und Hochladen von Code auf ein Arduino-Board.
	
	\item \textbf{Board wählen} Erkannte Arduino-Boards werden hier automatisch angezeigt, inklusive der Port-Nummer.
	
	\item \textbf{Sketchbook} Dieser Bereich dient als zentrale Ablage für alle Sketches, die lokal auf dem Computer gespeichert sind. Das Sketchbook bietet eine strukturierte, leicht zugängliche Umgebung zur Verwaltung, Organisation und Bearbeitung von Arduino-Code. Zusätzlich ermöglicht es die Synchronisation mit der Arduino Cloud, wodurch Sketches geräteübergreifend verfügbar sind. Diese Cloud-Integration erlaubt den Zugriff und die Bearbeitung von Sketches aus der Online-Arduino-Umgebung.
	
	\item \textbf{Board-Verwaltung} Durchsuchen von Arduino- und Drittanbieter-Paketen zur Installation. Beispielsweise erfordert ein MKR WiFi 1010 Board das Arduino SAMD Boards-Paket.
	
	\item \textbf{Bibliotheksverwalter} Zugriff auf tausende Arduino-Bibliotheken, bereitgestellt von Arduino und der Community.
	
	\item \textbf{Debuggen} Echtzeit-Test und Fehlersuche in Programmen.
	
	\item \textbf{Suchen} Durchsucht den geschriebenen Code nach Schlüsselwörtern.
	
	\item \textbf{Serieller Monitor} Öffnet den seriellen Monitor als neuen Tab in der Konsole.
	
	\item \textbf{Serieller Plotter} Visualisiert serielle Daten als Echtzeit-Diagramme.
	
\end{itemize}

\begin{figure}[H]
	\centering
	\caption{Die neue Arduino IDE mit der nützlichen Sidebar}
	\label{fig:ArduinoSidebar}
	\includegraphics[width=\textwidth]{Arduino/ArduinoIDE/arduinoSidebar.png}
\end{figure}


\section{Funktionen}
Die Arduino IDE 2 ist ein vielseitiges Entwicklungswerkzeug, das Funktionen wie die direkte Installation von Bibliotheken, Cloud-Synchronisation für Sketches und integrierte Debugging-Tools bietet. Dieser Abschnitt beleuchtet einige der Kernfunktionen, mit Verweisen auf detailliertere Ressourcen.

\subsection{Sketchbook}
Das \textbf{Sketchbook} ist der Speicherort für Arduino-Codedateien, die als Sketches bezeichnet werden. Diese Dateien werden mit der Endung \FILE{.ino} gespeichert. Zur ordnungsgemäßen Organisation muss jeder Sketch in einem Ordner abgelegt werden, der genau denselben Namen trägt wie die Sketch-Datei. Beispielsweise muss ein Sketch mit dem Namen \FILE{MySketch.ino} in einem Ordner namens \FILE{MySketch} gespeichert werden. Diese Namenskonvention ist entscheidend, damit die Arduino IDE die Sketches korrekt identifizieren und verwalten kann.

Standardmäßig werden Sketches in einem Ordner namens \FILE{Arduino} gespeichert, der sich im Dokumentenverzeichnis des Systems befindet.

Um auf das Sketchbook zuzugreifen, können Sie das Ordnersymbol in der Seitenleiste der Arduino IDE auswählen. Diese Aktion öffnet das Verzeichnis, in dem die Sketches gespeichert sind, und ermöglicht so eine effiziente Verwaltung und den Zugriff auf die Sketches. \ref{fig:arduinoSketch}

\begin{figure}[H]
	\centering
	\caption{Das Sketchbook auf der Arduino IDE}
	\label{fig:arduinoSketch}
	\includegraphics[width=\textwidth]{Arduino/ArduinoIDE/arduinoSketchbook.png}
\end{figure}


\subsection{Board-Verwaltung}
Die Board-Verwaltung kann unter folgendem Pfad gefunden werden: \menu{Werkzeuge > Board > Board-Verwaltung...}.

Der \textbf{Boards Manager} in der Arduino IDE ermöglicht die Installation verschiedener Board-Pakete, wie in Abbildung~\ref{fig:arduinoBoard} dargestellt. Ein Board-Paket enthält die erforderlichen Dateien und Anleitungen, um Code für bestimmte Arduino-Board-Typen zu kompilieren und hochzuladen.

Es stehen verschiedene Board-Pakete zur Verfügung, darunter:
\begin{itemize}
	\item \textbf{avr} - für ältere Boards wie das Arduino Uno
	\item \textbf{samd} - für neuere Boards wie das Arduino Zero
	\item \textbf{megaavr} - für die Mega-Serie
\end{itemize}

Jedes Paket unterstützt unterschiedliche Arduino-Board-Familien. Der Boards Manager stellt sicher, dass die richtigen Tools und Dateien für die Programmierung des ausgewählten Boards installiert sind.

\begin{figure}[H]
	\centering
	\caption{Die Board-Verwaltung auf der Arduino IDE}
	\label{fig:arduinoBoard}
	\includegraphics[width=\textwidth]{Arduino/ArduinoIDE/arduinoBoard.png}
\end{figure}


\subsection{Bibliotheksverwalter}
Den \textbf{Bibliotheksverwalter} erreichen Sie über folgenden Pfad:\\
 \menu{Werkzeuge > Bibliotheken verwalten...}

Der Bibliotheksverwalter in der Arduino IDE ermöglicht das Durchsuchen und Installieren einer Vielzahl von Bibliotheken. Diese Bibliotheken erweitern die Grundfunktionen von Arduino und vereinfachen die Arbeit mit speziellen Hardwarekomponenten, wie zum Beispiel:
\begin{itemize}
	\item Steuerung von Servomotoren
	\item Auslesen verschiedener Sensoren
	\item Kommunikation mit Modulen (z.B. Wi-Fi)
\end{itemize}

Die Bibliotheken enthalten vorgefertigten Code zur Ansteuerung spezifischer Hardware und vereinfachen komplexe Aufgabenstellungen. Dies ermöglicht eine schnellere und effizientere Entwicklung von Arduino-Projekten, wie in Abbildung~\ref{fig:arduinoBibliothek} dargestellt.

\begin{figure}[H]
	\centering
	\caption{Der Bibliotheksverwalter auf der Arduino IDE}
	\label{fig:arduinoBibliothek}
	\includegraphics[width=\textwidth]{Arduino/ArduinoIDE/arduinoBibliothek.png}
\end{figure}


\section{Integrating and Using a Custom Library in Arduino IDE}

In the development of embedded systems using the Arduino platform, libraries are essential for abstracting and simplifying complex functionality. A custom library is particularly beneficial when a specific functionality or hardware interface is repeatedly used across different projects. This section outlines the process of creating, installing, and using a custom library within an Arduino project, in the Arduino Integrated Development Environment (IDE).


\subsection{Creating the Library Files} 


	


A well-structured custom Arduino library typically consists of at least two key components:

	\menu{Header File (.h)}: This file contains the declarations of classes, functions, and variables that define the interface of the library. 
		It provides the necessary definitions for the functions and classes to be used by external programs.
		
		
	\menu{Source File (.cpp)}: This file contains the actual implementation of the functions and methods declared in the header file. 
		It defines the behavior of the functions and 	classes.
		The library structure is organized in a way that allows for clear separation between the declaration and implementation of its functionalities.






\subsection{Installation and Integration of the Library}

Once the custom library has been created, it must be properly installed and integrated into the Arduino IDE to be used in projects.

To install a custom library in Arduino follow these steps:

\begin{enumerate}
	\item Locate the Libraries Folder: The Arduino IDE searches for libraries in the libraries folder, which is located within the Arduino sketchbook directory. The typical path to this directory is:
	
	\medskip
	
	{\small\SHELL{ C:/Users/<username>/Documents/Arduino/libraries}}
	
	\item Copy the custom library folder, e.g., \FILE{Nano33BLESenseLE}, into the  directory \FILE{libraries}. The final path should look like this:
	
	\medskip
	
	{\small\SHELL{C:/Users/<username>/Documents/Arduino/libraries/Nano33BLESenseLED}}
	
	\item Restart the Arduino IDE: After placing the library in the correct directory, restart the Arduino IDE to ensure that it recognizes the newly added library.
	
\end{enumerate}





\subsection{Using Symbolic Links for Library Integration}

In some cases, it may be more convenient or necessary to store the library files outside the default library directory. For example, if the libraries are stored in a GitHub repository or for better version control, the command \menu{mklink} can be used to create a symbolic link. This allows the Arduino IDE to access the library files as if they were in the default directory, without duplicating the files.


\subsection{Windows Tool \SHELL{mklink} to Create Symbolic Links}

Symbolic links allow libraries to be stored in a different location while still being treated by the Arduino IDE as if they reside in the standard directory \FILE{libraries}. This can be particularly useful for version-controlled environments, such as when using Git, or to organize libraries without duplicating files.

Steps for Creating a Symbolic Link:



\begin{enumerate}



\item \textbf{Open Command Prompt with Administrator Privileges:} Press \menu{Win + S} and search for \menu{cmd}, click on \menu{Command Prompt} and select \menu{Run as Administrator} as seen in Figure \ref{fig:cmdLink}

\begin{center}\centering
	\includegraphics[width=8cm]{Arduino/ArdiunoIDE2/mklinkCmd.png}
	\captionof{figure}{Command Prompt with Administrator Privileges}
	\label{fig:cmdLink}		
\end{center}


\Mynote{Figure 1.17, 1.18 with an english sytem, so the figures are on english}



\item \textbf{Execute the Command mklink :} The command for creating a symbolic link is: {\small \SHELL {mklink /D "TargetPath" "SourcePath"}}

\begin{itemize}
	\item 	\textbf{TargetPath:} The location where the Arduino IDE expects the library to be, here \FILE{libraries}.
	\item 	\textbf{SourcePath:} The actual location of the library.
\end{itemize}
		
			
\end{enumerate}
		
\bigskip

		
For this example, the library is located at: 


{\small\SHELL{C:/Users/MSRLabor/Documents/GitHub/241031ArduinoNano33BLESense/report/Code/Nano33BLESense/Nano33BLESenseLED}}

\medskip


The target path in the Arduino IDE’s libraries folder is: 


{\small\SHELL{C:/Users/MSRLabor/Documents/Arduino/libraries/Nano33BLESenseLED}}

\medskip

The full command can be seen in Figure \ref{fig:cmdSynthax}.

\begin{center}\centering
	\includegraphics[width=10cm]{Arduino/ArdiunoIDE2/cmdSynthax.png}
	\captionof{figure}{Command Prompt}
	\label{fig:cmdSynthax}		
\end{center}


\textbf{Verifying the Symbolic Link:} After executing the command, navigate to the folder \FILE{libraries}. A folder named \FILE{Nano33BLESenseLED} should be present, acting as a symbolic link pointing to the actual library location.









\subsection{MacOS Command \SHELL{ln -s} to Create Symbolic Links}


\begin{enumerate}
	
	\item Open Terminal: Open the \menu{Terminal} application by searching for it in \menu{Applications}, go to \menu{Utilities}.
	
	\item Execute the Command \SHELL{ln -s} : The command for creating a symbolic link is: 
	
	\SHELL{ln -s "SourcePath" "TargetPath"}
	
	
	
	
	\begin{itemize}
		\item 	\textbf{TargetPath:} The location where the Arduino IDE expects the library to be, here \FILE{libraries}.
		\item 	\textbf{SourcePath:} The actual location of the library.
	\end{itemize}
	
\end{enumerate}
		
		
		
\bigskip

		
For this example, the library is located at: 


{\small\SHELL{C:/Users/MSRLabor/Documents/GitHub/241031ArduinoNano33BLESense/report/Code/Nano33BLESense/Nano33BLESenseLED}}

\medskip



The target path in the Arduino IDE’s libraries folder is: 


{\small\SHELL{C:/Users/MSRLabor/Documents/Arduino/libraries/Nano33BLESenseLED}}

\medskip

The full command would be: 

{\small\SHELL{ln -s "/Users/username/Documents/GitHub/241031ArduinoNano33BLESense/report/Code/Nano33BLESense/Nano33BLESenseLED" "/Users/username/Documents/Arduino/libraries/Nano33BLESenseLED"}}


\medskip


Verifying the Symbolic Link: After executing the command, navigate to the folder \FILE{libraries}. A folder named \FILE{Nano33BLESenseLED} should be present, acting as a symbolic link pointing to the actual library location.

\bigskip

\Mynote{here are some pictures necessary as in 1.5.4}





\bigskip



\subsection{Verifying Library Availability in the Arduino IDE}

To ensure that the custom library has been successfully integrated and is recognized by the Arduino IDE, follow these steps:

\begin{enumerate}
	
	\item Restart the Arduino IDE:
	If the Arduino IDE was open during the library installation or after creating the symbolic link, close and reopen it. This step ensures the IDE reloads its list of available libraries.
	
	\item Check the Library Menu: Navigate to \menu{Sketch} go to \menu{Include Library} and select the Library \FILE{Nano33BLESenseLED}, as seen in Figure \ref{fig:OwnLibrary}
	
	\item Verify Inclusion in the List:
	If the library is listed, it indicates successful recognition by the IDE. If not, double-check the directory structure and symbolic link to confirm they are correctly configured.
	
	\item Compile a Test Program:
	Write a simple test program using functions or classes from the library. Compile the sketch to ensure there are no errors, which confirms that the library has been successfully integrated.
	
\end{enumerate}


\begin{center}\centering
	\includegraphics[width=10cm]{Arduino/ArdiunoIDE2/includelibrary.png}
	\captionof{figure}{Including the Nano33BLESenseLED Library}
	\label{fig:OwnLibrary}		
\end{center}

Once the program compiles without errors, the library is ready to use, and you can confidently include it in your Arduino projects.












\subsection{Serial Monitor}


The Serial Monitor can be found in the following steps: \menu {Tool > Serial Monitor}

The Serial Monitor is a tool that allows to view data streaming from the board, via for example the command \PYTHON{Serial.print()}.

Historically, this tool was located in a separate window but is now integrated with the editor, making it easy to run multiple instances simultaneously on your computer. The Serial Monitor can be seen in Figure  ~\ref{serialMonitor}.


\begin{center}
        \includegraphics[width=0.7\linewidth]{Arduino/ArdiunoIDE2/serialMonitor.png}
        \captionof{figure}{Serial Monitor}\label{serialMonitor}
\end{center}


\subsection{Serial Plotter}

The tool Serial Plotter is great for visualizing data with graphs and monitoring, such as voltage peaks. It can monitor multiple variables simultaneously, with the option to enable specific types.


\section{Examples}

A significant component of the Arduino Documentation is the set of example sketches bundled with various libraries. These examples demonstrate the practical application of library functions, showcasing their intended uses and main features. Libraries included with certain board packages may also contain their own example sketches.
To access these example sketches whether from libraries installed manually or those included with board packages navigate to \menu{File > Examples}, and locate the desired library from the list.

For instance, when an UNO R4 WiFi board is connected, the examples list appears with options specific to that board. An example sketch demonstrating a pre-loaded Tetris animation can be accessed by navigating to \menu{File > Examples > LED Matrix > MatrixIntro} and uploading it to the connected board. This example shows how the board’s bundled code can be used in practice, assisting users in learning how to control various hardware functions directly. 

\begin{center}
        \includegraphics[width=0.7\linewidth]{Arduino/ArdiunoIDE2/examplesketches.png}
        \captionof{figure}{Examples}\label{examplesketches}
\end{center}


In Figure  ~\ref{examplesketches} above, the examples list is shown when a UNO R4 WiFi board is connected to the computer.



\section{Debugging}

The tool Debugger in Arduino IDE 2.3.3 is designed to assist in testing and debugging programs by providing the ability to step through code execution in a controlled manner. This tool allows users to set breakpoints, step through code line-by-line, and inspect variables, helping to identify issues such as logic errors or incorrect variable values. The debugger enables users to pause execution at specific points in the program, allowing for a detailed examination of the program's behavior and memory state. This feature enhances the development process by making it easier to locate and fix errors during runtime, rather than relying solely on print statements or trial and error. ~\ref{Debug}


\begin{center}
        \includegraphics[width=0.7\linewidth]{Arduino/ArdiunoIDE2/Debug.png}
        \captionof{figure}{Debugging Example}\label{Debug}
\end{center}



\Mynote{Debugging Bild muss ersetzt werden, weil nicht zu lesen}


\section{Autocompletion}

Autocompletion is a key feature in Arduino IDE version 2, making it easier to code by suggesting functions and elements from the Arduino API. This feature helps identify and complete code more quickly, improving efficiency and accuracy.

For autocompletion to work, the board must be selected in the IDE. This ensures that the editor recognizes the correct functions and libraries for the specific board, allowing accurate suggestions. ~\ref{autocomplete}


\begin{center}
        \includegraphics[width=0.7\linewidth]{Arduino/ArdiunoIDE2/autocomplete.png}
        \captionof{figure}{Autocomplete}\label{autocomplete}
\end{center}


\section{Conclusion}
This guide provides an overview of key features in Arduino IDE 2, along with references to detailed articles for further exploration. Each section aims to enhance understanding and enjoyment of the wide range of functionalities included in the IDE.



\begin{figure}[H]\centering
    \includegraphics[width=10cm]{Arduino/ArdiunoIDE2/MenuBarOption}
    \caption{Menu Bar Option}
    \label{fig::MenuBarOption}		
\end{figure}


%
%\chapter{To-Do}


%\begin{itemize}


%	 \item \textcolor{red}{Mac Instalation nochmal durchführen und aktualisieren (gilt auch für die Bilder)}
	 
	 
	 
%	\item \textcolor{red}{In Discription: Item-Punkte anpassen mit den Bildsymbolen} 
%	\item \textcolor{red}{MyNote Kommentare beachten}
%\end{itemize}

%






